\documentclass[11pt,a4paper]{article}
\usepackage[margin=27mm]{geometry}
\usepackage[T1]{fontenc}
\usepackage{lmodern,microtype,amsmath,amssymb,amsthm,mathtools,booktabs}
\usepackage[unicode,colorlinks=true,linkcolor=blue,citecolor=blue,urlcolor=blue]{hyperref}
\usepackage{enumitem}
\setlist{itemsep=3pt,topsep=5pt}
\newtheorem{theorem}{Theorem}[section]
\newtheorem{lemma}[theorem]{Lemma}
\newtheorem{proposition}[theorem]{Proposition}
\newtheorem{corollary}[theorem]{Corollary}
\theoremstyle{remark}\newtheorem{remark}[theorem]{Remark}
\newcommand{\D}{\mathbb D}
\newcommand{\R}{\mathbb R}
\newcommand{\C}{\mathbb C}
\newcommand{\av}[1]{\left\langle #1\right\rangle}
\newcommand{\Rea}{\operatorname{Re}}
\newcommand{\Ima}{\operatorname{Im}}
\newcommand{\Res}{\operatorname{Res}}
\newcommand{\supp}{\operatorname{supp}}
\newcommand{\norm}[1]{\left\lVert #1\right\rVert}
\DeclareMathOperator{\diag}{diag}
\title{Finite convex interpolation for polynomial Lyapunov existence}
\author{Lluis Eriksson}
\date{18 October 2026}
\begin{document}
\maketitle
\begin{abstract}
We characterize polynomial Lyapunov existence over all degrees for planar
homogeneous polynomial fields with nonvanishing angular velocity.
A moment threshold separates strict polynomial certificates from the
absence of local analytic weak certificates in the stable case.
Classical confluent Nevanlinna--Pick interpolation gives a finite convex
semialgebraic formulation, with exact boundary decision for algebraic
coefficients. For each odd field degree, the certificate region is open
and semialgebraic, and compact admissible families have a computable
common certificate degree.
For rational fields in dense binary encoding, we prove polynomial-time
additive approximation of the threshold and polynomial-time decision
under a supplied positive gap promise. Explicit inner and outer radii
permit a classical weak convex optimization reduction; a uniform
functional perturbation argument handles algebraic poles without
constructing their common number field. This approximates a value and
does not bound the output size of a Lyapunov polynomial polynomially.
Jackson bounds and a constructive sufficient-degree estimate complement
the finite formulation. An exact cubic example has an algebraic threshold
and minimum Lyapunov degree eight at one rational parameter.
Beyond the rotating class, a boundary-factor obstruction produces a
full twelve-dimensional rational box of stable quintic fields with
sign-changing angular velocity and no local analytic weak certificate.
One explicit rational $C^1$ Lyapunov function works for the entire box.
The latter result is a sufficient obstruction, not a general classification
with angular zeros. General certified solver implementations are not supplied.
\end{abstract}

\section{The certificate problem}
Consider a real planar homogeneous polynomial vector field \(X\) of
degree \(d\geq1\). Write its polar equations as
\begin{equation}\label{eq:polar}
 \dot r=r^d F(\theta),\qquad
 \dot\theta=r^{d-1}G(\theta).
\end{equation}
We assume that \(G\) never vanishes. Its sign is constant; reflecting
the plane if necessary makes \(G>0\). Antipodal parity then forces
\(d\) to be odd, and \(F,G\) are \(\pi\)-periodic trigonometric
polynomials. Define
\begin{equation}\label{eq:profile}
 a=-\frac1{2\pi}\int_0^{2\pi}\frac{F(\theta)}{G(\theta)}\,d\theta,
 \qquad
 q(\phi)=\frac{F(\phi/2)}{G(\phi/2)}+a.
\end{equation}
The profile \(q\) is real analytic and has zero mean.

A \emph{weak Lyapunov function} here satisfies \(V(0)=0\), is positive
away from the origin in its domain, and has \(L_XV\leq0\).
A strict function has \(L_XV<0\). Homogeneous polynomial certificates
are required globally; analytic certificates may be defined only near
the origin. We study the stable rotating case \(a>0\).

\begin{theorem}[Finite decision in the rotating class]\label{thm:decision}
For a planar homogeneous polynomial field with real algebraic
coefficients and \(G\) of one strict sign, existence of a strict
polynomial Lyapunov function is decidable by a finite algorithm.
In the case \(a>0\), the following are equivalent:
\begin{enumerate}[label=(\roman*)]
\item a strict homogeneous polynomial Lyapunov function exists;
\item a weak polynomial Lyapunov function exists;
\item a positive local real-analytic weak Lyapunov function exists.
\end{enumerate}
The algorithm also decides inputs on the boundary between existence
and nonexistence. It uses finite semialgebraic interpolation data,
without prescribing a Lyapunov degree. Exact boundary decision uses
quantifier elimination. Polynomial-time approximation and promised-gap
decision for dense rational inputs are established separately in
Theorem~\ref{thm:rationalcomplexity}.
\end{theorem}

Theorem~\ref{thm:decision} concerns existence of a particular kind of
certificate. Stability in this class already follows from the sign
of the classical polar average \(a\). The distinction matters:
Liu and Fitzsimmons \cite{liu} exhibit a stable homogeneous cubic
field with no polynomial or local analytic weak Lyapunov function.
Their family, stability analysis and original counterexample are
prior results. We use the same factorization and leading-term
viewpoint, extending the certificate problem to every profile in
\eqref{eq:profile}.

Earlier results show that polynomial Lyapunov degree cannot be
bounded solely by dimension and vector-field degree
\cite{ahmadi-parrilo,ahmadi-khadir}; this does not rule out an
input-dependent decision procedure. Fixed-degree semialgebraic
searches and converse sums-of-squares results address different
questions. Gram and Toeplitz methods also provide sufficient
certificates for systems with angular variables \cite{tripathi}.
Our finite interpolation formulation decides existence over all
polynomial degrees for the rotating homogeneous class.

Korda's recent stability decision procedure covers planar homogeneous
fields, including angular zeros \cite{korda}; Part~III of its supplement
explicitly distinguishes stability from polynomial-certificate existence.
Its Theorem~20.3 already computes the normalized polar mean as an
algebraic residue sum. We use that established calculation in our
normalization. The additional problem here is the existence of a
polynomial certificate of some degree.

The analytic tools used below are classical: Fej\'er--Riesz
factorization \cite{fejer}, positive-real representations,
Jackson kernels \cite{jackson}, semidefinite duality, confluent
interpolation \cite{clark}, and real quantifier elimination
\cite{tarski}. The contribution is the reduction connecting them
to this all-degree certificate problem.

\section{A moment threshold and finite error bounds}
For a real \(2\pi\)-periodic function use the conventions
\[
 \av{v}=\frac1{2\pi}\int_0^{2\pi}v(t)\,dt,\qquad
 q(t)=\sum_{k\in\mathbb Z}q_k e^{ikt},\qquad
 H\cos kt=\sin kt,\quad H\sin kt=-\cos kt.
\]
Assume for now that \(q\in C^2\), \(q_0=0\).
For \(N\geq1\), set
\[
 T_N(m)=(m_{k-j})_{j,k=0}^N,\qquad m_{-k}=\overline{m_k}.
\]
Define the finite optimum
\begin{equation}\label{eq:sdp}
 \begin{split}
 a_N(q)=\min\{b:\;&m_0=1,\quad T_N(m)\succeq0,\quad T_N(w)\succeq0,\\
 &w_0=b,\quad w_k=i m_k-q_k\quad(1\leq k\leq N)\}.
 \end{split}
\end{equation}
The diagonal of \(T_N(w)\) imposes \(b\geq0\).
The choice \(m_k=0\), \(b\geq\max q\), is feasible, since
\(b-q\geq0\). Principal minors give \(|m_k|\leq1\), so the
minimum is attained after restricting \(b\leq\max q\).
Principal submatrices also show
\[
 0\leq a_1(q)\leq a_2(q)\leq\cdots\leq\max q.
\]
Thus
\begin{equation}\label{eq:A}
 A(q)=\sup_{N\geq1}a_N(q)
\end{equation}
is well defined. The integer \(N\) is a moment order, not a
Lyapunov polynomial degree.

\begin{theorem}[Finite brackets]\label{thm:brackets}
Let \(\vartheta_N=\pi/(N+2)\). Then
\begin{equation}\label{eq:En}
 0\leq A(q)-a_N(q)
 \leq E_N:=\frac{\pi^2}{4}(1-\cos\vartheta_N)\norm{q''}_\infty
 \leq\frac{\pi^4\norm{q''}_\infty}{8(N+2)^2}.
\end{equation}
Independently, if the weighted Fourier sum converges,
\begin{equation}\label{eq:fourierbound}
 0\leq A(q)-a_N(q)
 \leq \frac{\pi^2}{(N+2)^2}\sum_{k\geq1}k^2|q_k|.
\end{equation}
The two right-hand majorants need not be ordered.
\end{theorem}

\begin{proof}
Put
\[
 v_j=\sqrt{\frac2{N+2}}\sin((j+1)\vartheta_N),\quad 0\leq j\leq N,
 \qquad K_N(t)=\left|\sum_{j=0}^N v_j e^{ijt}\right|^2.
\]
The kernel is even, nonnegative and has mean one. Its real Fourier
coefficients are
\[
 \kappa_k=\sum_{j=0}^{N-k}v_jv_{j+k}\ (0\leq k\leq N),
 \qquad \kappa_k=0\ (k>N),
\]
with \(\kappa_0=1\) and \(\kappa_1=\cos\vartheta_N\).
These are the classical Jackson coefficients:
\begin{equation}\label{eq:jackson}
 \kappa_k=
 \frac{(N+2-k)\cos(k\vartheta_N)+
       \cot(\vartheta_N)\sin(k\vartheta_N)}{N+2}
 \quad(0\leq k\leq N);
\end{equation}
see \cite[eq.~(2.11)]{jackson}.
Extending \(v\) by zero to \(\ell^2(\mathbb Z)\) and using the bilateral
shift gives
\begin{equation}\label{eq:shift}
 1-\kappa_k=\tfrac12\norm{v-S^kv}^2
 \leq k^2(1-\kappa_1),\qquad k\geq1.
\end{equation}

Take an optimizer in \eqref{eq:sdp} and set
\[
 u_N(t)=1+2\Rea\sum_{k=1}^N\kappa_k m_k e^{ikt},\qquad
 w_N(t)=a_N+2\Rea\sum_{k=1}^N\kappa_k w_k e^{ikt}.
\]
For \(z_j(t)=v_j e^{ijt}\), these are respectively
\(z(t)^*T_N(m)z(t)\) and \(z(t)^*T_N(w)z(t)\).
Both are nonnegative, and
\[
 w_N=a_N-q_N-Hu_N,\qquad q_N=K_N*q.
\]
With \(E=\norm{q-q_N}_\infty\), the two nonnegative densities
\(u_N\) and \(w_N+E-(q-q_N)\) give feasible moments at every
order at damping \(a_N+E\). Hence \(A\leq a_N+E\). This argument
constructs compatible full moment sequences; it does not assume
that arbitrary simultaneous truncations extend.

Taylor's theorem, evenness of the kernel, and
\(t^2\leq(\pi^2/2)(1-\cos t)\) on \([-\pi,\pi]\) yield
\[
 E\leq\frac{\norm{q''}_\infty}{2}\av{t^2K_N(t)}
 \leq\frac{\pi^2}{4}\norm{q''}_\infty(1-\kappa_1).
\]
Finally \(1-\cos\vartheta_N\leq\vartheta_N^2/2\) gives
\eqref{eq:En}. Applying \eqref{eq:shift} directly to the Fourier
series gives
\(E\leq2(1-\kappa_1)\sum_{k\geq1}k^2|q_k|\), proving
\eqref{eq:fourierbound}.
\end{proof}

\section{All degrees of polynomial certificates}
For \(|\alpha|<1\), define the positive quadratic form
\begin{equation}\label{eq:quad}
 Q_\alpha(x,y)=(1+|\alpha|^2)(x^2+y^2)
 -2\Rea(\alpha)(x^2-y^2)+4\Ima(\alpha)xy.
\end{equation}
Its circle profile is \(|1-\alpha e^{2i\theta}|^2\).

\begin{theorem}[Strict threshold]\label{thm:threshold}
For \(a>0\) and a real mean-zero \(C^2\) profile \(q\), the following
are equivalent:
\begin{enumerate}[label=(\roman*)]
\item \(a>A(q)\);
\item a finite product \(P=\prod_{j=1}^M Q_{\alpha_j}\),
      \(|\alpha_j|<1\), satisfies
\[
  -a+q(2\theta)+\frac{p'(\theta)}{2M p(\theta)}<0,
  \qquad p(\theta)=P(\cos\theta,\sin\theta);
\]
\item some positive homogeneous polynomial of even degree satisfies
      the corresponding weak inequality with \(\leq0\).
\end{enumerate}
In particular equality \(a=A(q)>0\) admits no such weak certificate.
\end{theorem}

\begin{proof}
For a positive homogeneous polynomial of degree \(2M\),
Fej\'er--Riesz factorization of its circle profile, choosing the
factor with all zeros outside the closed disk, gives
\[
 p(\theta)=C\prod_{j=1}^M|1-\alpha_j e^{2i\theta}|^2,
 \qquad |\alpha_j|<1.
\]
Factors with \(\alpha_j=0\) pad the actual Fourier degree to \(M\).
Differentiation gives
\begin{equation}\label{eq:herglotz}
 \frac{p'(\theta)}{2M p(\theta)}
 =\Ima h(e^{2i\theta}),\qquad
 h(z)=\frac1M\sum_{j=1}^M\frac{1+\alpha_jz}{1-\alpha_jz},
 \quad h(0)=1.
\end{equation}
Thus \(u=\Rea h\) has a strictly positive minimum on the circle.
If \(w=a-q-Hu\geq0\), its coefficients satisfy
\(w_k=i m_k-q_k\), where \(m_k\) are the coefficients of \(u\).
Both moment matrices are positive, so \(A\leq a\).

To prove the necessary \emph{strict} inequality, let
\(Q(z)=2\sum_{k\geq1}q_kz^k\). The \(C^2\) hypothesis implies
absolute convergence on the closed disk. For small \(\varepsilon>0\),
\[
 h_\varepsilon=(1-\varepsilon)h+\varepsilon(1-iQ)
\]
still has positive real part and \(h_\varepsilon(0)=1\).
Since \(\Ima(1-iQ)=-q\) on the circle, its slack is
\[
 a-q-\Ima h_\varepsilon=(1-\varepsilon)w+\varepsilon a.
\]
It remains positive after replacing \(a\) by
\(a-\varepsilon a/2\). This gives \(A<a\), using \(a>0\).

Conversely suppose \(a>A\). Theorem~\ref{thm:brackets} supplies
\(N\) with \(a_N+E_N<a\). Form the nonnegative densities \(u_N,w_N\)
in its proof, and let \(h_N(z)=1+2\sum\kappa_km_kz^k\).
For \(r<1\) close enough to one,
\[
 \Rea h_N(rz)>0\quad(|z|\leq1),\qquad
 a-q-\Ima h_N(re^{it})
 =(w_N)_r+(a-a_N)-[q-(q_N)_r]>0.
\]
Here a subscript \(r\) means Poisson smoothing.
The Herglotz representation of \(h_N\) has a probability measure
on the circle. Approximate it by equal-weight empirical measures.
For this fixed \(r<1\), the kernel is uniformly continuous on the
closed disk, so the resulting functions
\[
 h_M(z)=\frac1M\sum_{j=1}^M
        \frac{1+r e^{-it_j}z}{1-r e^{-it_j}z}
\]
approximate \(h_N(rz)\) uniformly. They retain the strict slack
when \(M\) is large. Equation~\eqref{eq:herglotz} now gives the
required product of positive quadratic forms.
\end{proof}

\begin{proposition}[Local analytic and global polynomial certificates]
\label{prop:analytic}
For the field \eqref{eq:polar} with \(G>0\) and \(a>0\), the three
conditions in Theorem~\ref{thm:decision} are equivalent to \(a>A(q)\).
\end{proposition}
\begin{proof}
For \(P=r^{2M}p(\theta)\),
\begin{equation}\label{eq:lie}
 \frac{L_XP}{2M r^{d-1}G P}
 =-a+q(2\theta)+\frac{p'(\theta)}{2M p(\theta)}.
\end{equation}
Thus Theorem~\ref{thm:threshold} handles homogeneous certificates.
If \(V\) is local analytic, take its first nonzero homogeneous
Taylor term \(P_m\). Positivity of \(V\) and homogeneity of \(X\)
give \(P_m\geq0\) and \(L_XP_m\leq0\).
For \(p_m=P_m|_{S^1}\),
\[
 p_m'+m(F/G)p_m\leq0.
\]
If \(p_m(\theta_0)=0\), multiplication by
\(\exp(m\int_{\theta_0}^{\theta}F/G)\) yields a nonnegative
nonincreasing function starting at zero, over a full following
period. This forces \(p_m\equiv0\), a contradiction. Hence \(P_m\)
is positive definite, \(m\) is even, and
Theorem~\ref{thm:threshold} gives \(a>A(q)\).
A polynomial is analytic; a strict homogeneous polynomial is
itself a local analytic weak certificate.
\end{proof}

The local leading-term mechanism follows the argument used for the
specific cubic counterexample in \cite{liu}.
For completeness, the classical stability statement is also direct.
A periodic primitive \(B'(\theta)=F/G+a\) gives
\[
 V_{\rm sm}(r,\theta)=r^2 e^{-2B(\theta)},\qquad
 L_XV_{\rm sm}=-2a r^{d-1}G V_{\rm sm}.
\]
This function is coercive and \(C^1\) at the origin; for \(a>0\)
it proves global asymptotic stability. If \(a=0\) the polar trajectories
are closed, and if \(a<0\) they spiral outward. Neither case admits
a strict Lyapunov function. This does not exclude weak first integrals
when \(a=0\).

\section{A finite interpolation representation}
We now assume
\begin{equation}\label{eq:rational}
 q(t)=R(e^{it}),
\end{equation}
where \(R\) is rational, has no pole on the unit circle, is real
there and has zero circle mean. Algebraic coefficients will be
needed only for the effective decision conclusion.

\subsection{The quadrant dual}
\begin{lemma}[Finite dual]\label{lem:dual}
The number \(a_N(q)\) equals the maximum of
\begin{equation}\label{eq:dualobjective}
 \av{q\,\Ima f}-c
\end{equation}
over analytic polynomials \(f\) of degree at most \(N\) such that
\(f(0)=c+i\), \(c\geq0\), and \(\Rea f,\Ima f\geq0\) on the circle.
\end{lemma}
\begin{proof}
Positive semidefinite Gram multipliers for the two Toeplitz
constraints correspond to nonnegative trigonometric polynomials
\(U,V\) of degree at most \(N\). Every such polynomial has a Gram
representation by Fej\'er--Riesz. With formal moments of \(u\),
\(\av{u}=1\), the Lagrangian is
\[
 b-\av{uU}-\av{(b-q-Hu)V}
 =b(1-\av V)+\av{qV}-\av{u(U+HV)}.
\]
We used skew-adjointness of \(H\) under the real circle pairing.
Stationarity requires \(\av V=1\) and \(U+HV=c\) constant,
with \(c=\av U\geq0\). This says exactly that \(f=U+iV\) is
analytic, \(f(0)=c+i\). The dual objective is \eqref{eq:dualobjective}.
The primal is strictly feasible at \(m_k=0\), \(b>\max q\);
its finite optimum is attained. Finite-dimensional semidefinite
strong duality proves the assertion.
\end{proof}

Let \(\zeta_j\) be the interior poles of \(R(z)/z\), of orders
\(n_j\). For any analytic \(f\) on \(\D\), define the finite jet functional
\begin{equation}\label{eq:lambda}
 \Lambda_R(f)=\sum_{|\zeta_j|<1}
       \Res_{z=\zeta_j}\frac{R(z)f(z)}{z}.
\end{equation}
Only derivatives through order \(n_j-1\) at \(\zeta_j\) occur.
For polynomials, the residue theorem gives
\(\Lambda_R(f)=\av{qf}\). For a general analytic \(f\),
\eqref{eq:lambda} is the definition; a boundary integral is not required.

\begin{proposition}[Quadrant formulation]\label{prop:quadrant}
Let \(\mathcal Q_c\) be the analytic maps from \(\D\) to the open
first quadrant with \(f(0)=c+i\), where \(c>0\). Then
\begin{equation}\label{eq:quadrant}
 A(q)=\sup_{c>0,\ f\in\mathcal Q_c}
           \bigl(\Ima\Lambda_R(f)-c\bigr).
\end{equation}
\end{proposition}
\begin{proof}
Any dual polynomial with \(c>0\) belongs to \(\mathcal Q_c\)
by the strong minimum principle for its two harmonic parts.
A dual polynomial with \(c=0\) is identically \(i\) and contributes
zero. Constants \(c+i\) already have objectives tending to zero
as \(c\downarrow0\).

For \(f\in\mathcal Q_c\) and \(r<1\), both parts of \(f(rz)\)
have positive minima on the closed disk. Its Taylor polynomials
converge uniformly there and eventually preserve both positive
parts, with the same constant term. Their objectives approach
\(\Ima\Lambda_R(f(r\,\cdot))-c\). Every interior jet in
\eqref{eq:lambda} converges to the jet of \(f\) as \(r\uparrow1\).
The union of the finite duals therefore has the same supremum
as \eqref{eq:quadrant}.
\end{proof}

\subsection{Interior jets and a Pick matrix}
Include \(0\) among the nodes, adding multiplicity one if it is
not already a pole of \(R(z)/z\). Let the distinct nodes be
\(\zeta_1,\ldots,\zeta_s\), with multiplicities \(n_j\), and put
\(m=\sum n_j\). Introduce normalized jet variables
\[
 u_{jp}=\frac{f^{(p)}(\zeta_j)}{p!},\qquad 0\leq p<n_j.
\]
At node \(0\), impose the value \(c+i\). Require
\begin{equation}\label{eq:jetpositive}
 c>0,\qquad \Rea u_{j0}>0,\quad \Ima u_{j0}>0.
\end{equation}
Define the jets of the proposed square \(g=f^2\) by
\[
 v_{jp}=\sum_{\ell=0}^p u_{j\ell}u_{j,p-\ell}.
\]
Form the Hermitian \(m\times m\) matrix
\begin{equation}\label{eq:pick}
 \mathcal P_{(j,p),(k,\ell)}
 =\left.\frac{\partial_z^p\partial_{\bar w}^{\,\ell}}{p!\ell!}
 \frac{g(z)-\overline{g(w)}}{i(1-z\bar w)}
 \right|_{z=\zeta_j,\ w=\zeta_k}.
\end{equation}
The right-hand side uses only the specified jets \(v\).
In particular for values alone its entry is
\((v_{j0}-\overline{v_{k0}})/(i(1-\zeta_j\bar\zeta_k))\).

\begin{lemma}[Classical confluent interpolation]\label{lem:pick}
The tuples satisfying \eqref{eq:jetpositive}, the value condition at
\(0\), and \(\mathcal P\succeq0\), are exactly the jet tuples of
the functions in \(\bigcup_{c>0}\mathcal Q_c\).
\end{lemma}
\begin{proof}
If \(f\) maps to the first quadrant, \(g=f^2\) maps to the
upper half-plane. Its positive Pick kernel and its derivatives give
\(\mathcal P\succeq0\).

Conversely, classical interior confluent Nevanlinna--Pick
interpolation provides an analytic \(g\) with \(\Ima g\geq0\)
and the prescribed jets when \(\mathcal P\succeq0\), including
singular semidefinite data. This is the half-plane theorem;
see \cite[Section~3, pp.~246--247]{clark}, after multiplication
of \(g\) by \(-i\). The normalized derivative kernels in
\eqref{eq:pick} and the repeated-node basis there are related by
an invertible change of basis. All interpolation nodes are interior.

At zero, \(g(0)=(c+i)^2\) has imaginary part \(2c>0\);
the minimum principle implies \(\Ima g>0\) throughout \(\D\).
The unique analytic square root taking values in the first
quadrant exists on the upper half-plane. The inequalities in
\eqref{eq:jetpositive} select this root at each node.
Successive derivatives of \(f^2=g\) recover all higher prescribed
jets, since the coefficient of the next jet is \(2u_{j0}\ne0\).
\end{proof}

\begin{theorem}[Finite algebraic threshold]\label{thm:algebraic}
For a rational profile \eqref{eq:rational}, \(A(q)\) is the
supremum of the finite problem
\begin{equation}\label{eq:finite}
 \sup\{J(c,u): (c,u)\in\mathcal S_R\},\qquad
 J(c,u)=\Ima\Lambda_R(u)-c,
\end{equation}
where \(\mathcal S_R\) is defined by \eqref{eq:jetpositive}, the
value at \(0\), and \(\mathcal P\succeq0\).
If \(R\) has algebraic coefficients, \(\mathcal S_R\) is
semialgebraic over real algebraic numbers, \(A(q)\) is algebraic,
and its comparison with any algebraic number is decidable.
The supremum in \eqref{eq:finite} need not be attained.
\end{theorem}
\begin{proof}
Proposition~\ref{prop:quadrant} and Lemma~\ref{lem:pick} give
\eqref{eq:finite}. Interior poles and residue coefficients of
an algebraic rational function are algebraic and effectively
computable. Each matrix entry in \eqref{eq:pick} is a polynomial
of degree at most two in the real jet variables, with algebraic
coefficients. The denominators \(1-\zeta_j\bar\zeta_k\) are
nonzero constants. Hermitian positivity can be expressed using
principal minors, or a real symmetric block matrix.
Thus \(\mathcal S_R\) is semialgebraic and \(J\) is affine.

The set is nonempty, and its supremum is finite and nonnegative
by \eqref{eq:A}. For algebraic \(a>0\), the condition \(a>A(q)\)
is the first-order real sentence
\begin{equation}\label{eq:qe}
 \exists\varepsilon>0\ \forall(c,u)\quad
 [(c,u)\in\mathcal S_R\ \Longrightarrow\ J(c,u)\leq a-\varepsilon].
\end{equation}
Real quantifier elimination \cite{tarski} decides this sentence.
The supremum itself is defined by the conditions that \(B\)
is an upper bound and that, for each \(\varepsilon>0\), some
feasible \(J\) exceeds \(B-\varepsilon\).
Eliminating quantifiers describes the singleton \(\{A(q)\}\)
over the real algebraic numbers. Such a singleton is algebraic,
and real root isolation gives an effective specification.
No compactness or maximizer of \eqref{eq:finite} is assumed.
\end{proof}

\begin{proof}[Proof of Theorem~\ref{thm:decision}]
The no-zero condition on the angular component is decidable on
the unit circle by real quantifier elimination. Reflect if needed.
The functions \(F,G\) have only even angular frequencies, so \(F/G\)
is rational in \(z=e^{2i\theta}\) with algebraic coefficients and
no circle poles. Its mean is a sum of interior residues of
\((F/G)(z)/z\), and hence \(a\) in \eqref{eq:profile} is algebraic
and computable; see also \cite[Supplement, Theorem~20.3]{korda}.
For \(a\leq0\) there is no strict certificate,
as noted after Proposition~\ref{prop:analytic}. For \(a>0\),
form \(R=F/G+a\), its interior nodes and the finite problem
\eqref{eq:finite}. Decide \eqref{eq:qe} and apply
Proposition~\ref{prop:analytic}. Equality \(a=A(q)>0\) yields
a negative answer by the strict threshold theorem.
\end{proof}

\begin{remark}
The matrix \(\mathcal P\) has size determined by the poles of the
input profile, rather than a candidate Lyapunov degree.
Its entries are quadratic in the jets, so this representation is not
a linear matrix inequality in the displayed variables. Nevertheless,
\(\mathcal S_R\) is convex: convex combinations of quadrant maps with
\(\Ima f(0)=1\) remain in that class, and taking jets is linear.
Lemma~\ref{lem:pick} identifies its image with \(\mathcal S_R\).
Thus \eqref{eq:finite} is a convex semialgebraic optimization problem.
This observation alone supplies neither a finite lifted semidefinite
representation nor an efficiency bound. The finite moment programs
\eqref{eq:sdp} are also convex, but their approximation alone does not
decide equality.
\end{remark}

\section{Polynomial-time additive approximation}
The convexity of the jet set has a quantitative consequence. The
input model throughout this section is dense binary rational data:
all coefficients or jets up to the indicated degree are explicit.
The statements do not apply to compressed exponent encodings.

\begin{theorem}[Rational jet data]\label{thm:ratjets}
Let \(\zeta_j\in\mathbb Q[i]\cap\D\) be distinct nodes including
zero, with positive multiplicities \(n_j\), and let
\[
 \Lambda(f)=\sum_{j,p<n_j}b_{jp}\frac{f^{(p)}(\zeta_j)}{p!},
 \qquad b_{jp}\in\mathbb Q[i],\quad\Lambda(1)=0.
\]
Over quadrant maps \(f(0)=c+i\), \(c>0\), define
\(A_\Lambda=\sup(\Ima\Lambda(f)-c)\).
For rational \(0<\varepsilon\leq1\), a rational approximation to
\(A_\Lambda\) with absolute error at most \(\varepsilon\) can be
computed in time polynomial in the binary encoding of the data and
the tolerance. In particular, the cost for \(\varepsilon=2^{-p}\)
is polynomial in \(p\). This is a value approximation, not a
polynomial-time construction of a Lyapunov polynomial.
\end{theorem}
\begin{proof}
Write \(f=c+iH\). Then \(H(0)=1\), \(\Rea H>0\), and
Herglotz representation gives, for \(|z|\leq r<1\),
\begin{equation}\label{eq:hergbounds}
 |H(z)|\leq U_0(r):=\frac{1+r}{1-r},\qquad
 \left|\frac{H^{(p)}(z)}{p!}\right|
 \leq U_p(r):=\frac2{(1-r)^{p+1}}\quad(p\geq1).
\end{equation}
Take the rational radii \(r_j=(1+|\zeta_j|^2)/2\), which dominate
\(|\zeta_j|\), and put
\[
 B=1+\sum_{j,p}(|\Rea b_{jp}|+|\Ima b_{jp}|)U_p(r_j).
\]
Since \(\Lambda(1)=0\), the affine objective is
\(J=\Rea\Lambda(H)-c\leq B-c\). Constants \(c+i\) show
\(0\leq A_\Lambda\leq B\), and near-maximizers may be chosen
with \(c\leq B+1\).

For \(0<\delta\leq1/2\), let \(K_\delta\) be the feasible jet
set with \(\delta\leq c\leq C:=2B+2\). It is convex. Harnack
gives the positive closed value bounds
\[
 \Rea f(\zeta_j)\geq\delta\frac{1-r_j}{1+r_j},\qquad
 \Ima f(\zeta_j)\geq\frac{1-r_j}{1+r_j}.
\]
Add the interval for \(c\) and the Pick PSD condition.
Lemma~\ref{lem:pick} shows that these closed conditions characterize
\(K_\delta\) exactly.
Equation~\eqref{eq:hergbounds} bounds all remaining coordinates,
so this body is compact. Take a nearly optimal map. Its convex
combination with the constant \(1+i\), using weights
\(1-\delta,\delta\), shows
\begin{equation}\label{eq:truncate}
 0\leq A_\Lambda-\max_{K_\delta}J\leq\delta(B+1).
\end{equation}
This does not assume attainment of the original supremum.

Put \(m=\sum n_j\) and remove the fixed coordinate
\(\Ima f(0)=1\), leaving \(n=2m-1\) real coordinates.
Let \(x_0\) be the jets of \(1+i\), and let the confluent
Vandermonde matrix be
\[
 V_{(j,p),\ell}=\binom\ell p\zeta_j^{\ell-p}
 \quad(0\leq\ell<m),
\]
with zero entries if \(\ell<p\). It is invertible. Define
\[
 T=\max\left(1,\sum_{a,b}
 (|\Rea(V^{-1})_{ab}|+|\Ima(V^{-1})_{ab}|)\right),
 \qquad r_0=(4T)^{-1}.
\]
Any real jet perturbation of Euclidean norm at most \(r_0\)
has a Hermite interpolating polynomial \(p\) with
\(\sup_{\overline\D}|p|\leq Tr_0\leq1/4\).
Thus \(1+i+p\) stays in the quadrant and its anchor has real
part in \([3/4,5/4]\). Consequently
\(B(x_0,r_0)\subseteq K_\delta\).
With \(U=\max_{j,p}U_p(r_j)\), the rational radius
\(R_0=2n(C+U+2)\) contains \(K_\delta\) about both zero and
\(x_0\).

Membership of rational tuples in \(K_\delta\) is polynomial-time:
form the quadratic squared-map jets and the rational Hermitian
Pick matrix, test the scalar bounds, and test PSD by exact Schur
elimination. A negative diagonal rejects; a zero diagonal requires
its row to be zero; a positive pivot allows a Schur complement.
Rational pivots are ratios of minors, so determinant bounds control
their bit lengths polynomially. No enumeration of all principal
minors is needed. This exact test is a valid weak-membership oracle.

Apply the classical centered-convex-body oracle theorem of
Gr\"otschel--Lov\'asz--Schrijver
\cite[Corollary~4.3.12]{gls}. Its hypotheses are the rational
center, inner and outer radii just exhibited and weak membership.
For tolerance \(\eta<r_0\), it returns \(\widehat x\) within
\(\eta\) of \(K_\delta\), whose objective is at least the
maximum on its \(\eta\)-interior minus \(\eta\).
The empty-interior alternative is impossible because of the inner ball.
Write \(J(x)=w\cdot x+j_0\), \(W=\max(1,\sum|w_i|)\).
Contracting a maximizer toward \(x_0\) by \(\eta/r_0\)
puts it in that interior and loses at most \(\eta W R_0/r_0\).
Proximity bounds any excess objective by \(\eta W\). Therefore
\[
 |J(\widehat x)-\max_{K_\delta}J|
 \leq\eta[1+W(1+R_0/r_0)].
\]
Choose
\[
 \delta=\frac\varepsilon{4(B+1)},\qquad
 \eta=\min\left(\frac{r_0}2,
 \frac\varepsilon{4[1+W(1+R_0/r_0)]}\right).
\]
Together with \eqref{eq:truncate}, this gives error at most
\(\varepsilon/2\). All displayed rational data have polynomial
binary length: powers have exponents \(O(m)\), and rational
matrix inversion and determinant bounds control the inverse
Vandermonde. Tiny geometric radii affect their logarithms, not
the number of bits exponentially. The oracle theorem bounds
both the number and size of its queries, completing the ordinary
Turing polynomial-time reduction.
\end{proof}

The theorem uses classical weak optimization, not a new general
convex optimization algorithm or an asserted finite LMI lift.
It already applies to trigonometric polynomials with rational Fourier
coefficients, since their functionals use only jets at zero.
Equality and the size of an explicit certificate remain separate issues.

\subsection{Rational profiles with algebraic poles}
Exact arithmetic in the compositum of all pole fields need not have
polynomial degree. The next argument avoids that arithmetic by
approximating the objective on a fixed analytic function class.

\begin{lemma}[Uniform perturbation of the functional]\label{lem:perturb}
Let \(f_p=f^{(p)}/p!\), and suppose that the functionals
\[
 \Lambda(f)=\sum b_{jp}f_p(\zeta_j),\qquad
 \widetilde\Lambda(f)=\sum\widetilde b_{jp}f_p(\widetilde\zeta_j)
\]
annihilate constants. Suppose corresponding
nodes have the same jet orders and all nodes have modulus at most
\(\rho<1\). Then
\[
 |A_\Lambda-A_{\widetilde\Lambda}|\leq
 \sum_{j,p}|b_{jp}-\widetilde b_{jp}|U_p(\rho)
 +\sum_{j,p}|\widetilde b_{jp}|(p+1)U_{p+1}(\rho)
 |\zeta_j-\widetilde\zeta_j|.
\]
\end{lemma}
\begin{proof}
Write \(f=c+iH\) as before. Constants cancel exactly in both
objectives. Since \(H_p'=(p+1)H_{p+1}\), integrate the derivative
bound along the chord between corresponding nodes. Equation~
\eqref{eq:hergbounds} bounds the coefficient error. Taking suprema
over the same class of quadrant maps proves the assertion.
No continuity claim about perturbed Pick feasible sets is used.
\end{proof}

\begin{theorem}[Polynomial-time value approximation for rational fields]
\label{thm:rationalcomplexity}
Let \(S=N/D\), \(N,D\in\mathbb Q[i][z]\), be given densely,
real on the unit circle, with \(D\) having no zeros there.
Both \(\mu=\av{S}\) and \(A(S-\mu)\) admit rational additive
\(\varepsilon\)-approximations in time polynomial in the dense
binary input length and the tolerance encoding.

Consequently, given a rational planar homogeneous field with
nonvanishing angular speed, and a rational promised separation
\(\gamma>0\), existence of a strict polynomial Lyapunov function
of some degree can be decided in polynomial time on inputs satisfying
\(|a-A(q)|\geq\gamma\). For stable inputs \(a>0\), this also
decides existence of a local analytic weak certificate under the
same promise. The latter equivalence is not asserted for centers.
\end{theorem}
\begin{proof}
Clear a common denominator of \(N,D\), scaling both together,
and put \(Q=zD\). At an interior root \(\zeta_j\) of \(Q\)
of multiplicity \(n_j\), set
\[
 E_j(t)=\sum_{k\geq0}\frac{Q^{(n_j+k)}(\zeta_j)}{(n_j+k)!}t^k,
 \qquad c_{jp}=[t^{n_j-1-p}]\frac{N(\zeta_j+t)}{E_j(t)}.
\]
Then \(\Lambda_S(f)=\sum c_{jp}f_p(\zeta_j)\) is the residue
functional, and \(\mu=\sum_jc_{j0}\). Include the anchor zero
with multiplicity at least one. The normalized functional
\(\Lambda=\Lambda_S-\mu f(0)\) has the same derivative
coefficients, and its zeroth anchor coefficient is
\[
 b_{00}=-\sum_{\zeta_j\ne0}c_{j0}.
\]
Thus it annihilates constants without first constructing the
algebraic mean in a common number field. Removable poles may be
cancelled by rational polynomial gcd, or retained with zero residues.

We record the bit bounds needed to approximate these data. Let
\(n,\tau\) bound degrees and cleared integer coefficient lengths;
both are polynomial in the dense input. Define
\(Q^\#(z)=z^{\deg Q}\overline{Q(1/\bar z)}\).
The primitive squarefree part \(P\) of
\[
 zQQ^\#\overline Q\,\overline{Q^\#}
\]
is an integer polynomial of polynomial degree and coefficient
length; bars on polynomial factors here conjugate coefficients.
Its roots include zero, all roots of \(Q\), and the reciprocal
conjugates of their nonzero roots. Squarefree extraction uses
polynomial gcd and preserves polynomial bit bounds.

For \(P\) of degree \(d_0\geq2\) and height \(H_0\geq1\),
Cauchy's root bound and its nonzero integral discriminant give the
conservative separation
\(\sigma=(4H_0)^{-d_0^2}\) between distinct roots. Indeed all
roots have modulus at most \(2H_0\), and isolating one squared
distance in the discriminant product, while bounding every other
distance by \(4H_0\), proves this lower bound. When \(d_0=1\)
only zero is relevant and we set \(\sigma=1\).
For an interior nonzero root with \(|\zeta|>1/2\),
\[
 \sigma\leq|\zeta-1/\bar\zeta|
 =\frac{1-|\zeta|^2}{|\zeta|},\qquad
 1-|\zeta|\geq\sigma/4.
\]
Hence \(\beta=\min(1/2,\sigma/4)\) is a lower bound for
the distance of every interior node from the unit circle.
The reciprocal argument likewise separates exterior roots from it.
Both \(\sigma\) and \(\beta\) have polynomial binary length.

The multiplicities and number of distinct roots of \(Q\) are
computed by gcd over \(\mathbb Q[i]\). Certified complex root
isolation and refinement have polynomial bit complexity here;
for example, the geometric bounds of
Mehlhorn--Sagraloff--Wang \cite[Theorems~3--4]{msw} apply.
Root magnitudes, inverse separations and inverse products of
root differences in those bounds have polynomial logarithms by
the preceding estimates. The rational coefficient oracle is exact.
Refining to \(2^{-k}\) accuracy therefore costs a polynomial in
the original input and \(k\).

Residue coefficients have the same approximation property, even
at multiple poles. At multiplicity \(r\),
\[
 \left|\frac{Q^{(r)}(\zeta)}{r!}\right|
 =|\operatorname{lead}(Q)|\prod_{w\ne\zeta}
 |\zeta-w|^{\operatorname{mult}(w)}\geq\sigma^{\deg Q}.
\]
The leading coefficient is a nonzero Gaussian integer. For an
explicit bound let \(n\geq1\) dominate both polynomial degrees,
let \(2^\tau\) bound real and imaginary coefficient parts, and put
\[
 C_0=(n+1)^3 2^{\tau+2n+4},\qquad
 h=\sigma^{\deg Q},\qquad B_0=2C_0/h.
\]
On \(|z|\leq2\), \(C_0\) bounds all factorial-normalized
derivatives of \(N,Q\) and their first derivatives, by the
binomial theorem. Within distance
\(\rho_0=\min(1/4,h/(2C_0))\) of the root, the series
denominator's constant term has modulus at least \(h/2\).
Write its coefficients as \(e_j(w)\), and those of the translated
numerator as \(a_j(w)\). Series division obeys
\[
 t_\ell=\frac{a_\ell-\sum_{j=1}^\ell e_jt_{\ell-j}}{e_0}.
\]
Induction gives \(|t_\ell|\leq B_0(1+B_0)^\ell\).
Differentiating the recurrence gives
\[
 |t'_\ell|\leq B_0\left(1+\sum_{k<\ell}|t_k|
                  +\sum_{k<\ell}|t'_k|+|t_\ell|\right).
\]
A second induction bounds every derivative for \(\ell\leq n\)
by \(L_0=(n+2)B_0^2(1+B_0)^{2n}\), whose logarithm is
polynomial in the input. Root accuracy
\(\min(\rho_0,\delta/(2L_0))\), followed by rounding within
\(\delta/2\), gives residue error at most \(\delta\).
Exact rational series evaluation has polynomial bit growth: at
order \(\ell\) its denominator introduces at most \(\ell+1\)
powers of the common constant denominator, in addition to
polynomial coefficient and evaluation denominators.
Thus the computation is polynomial in the input and requested
precision bits, treating each pole separately.

Let \(m\) be the total interior multiplicity including the anchor.
Coarse certified coefficient enclosures supply rational
\(M\geq1+\max|b_{jp}|\) with polynomial length. Set
\[
 \rho=1-\beta/2,\qquad U=\max_{0\leq p\leq m}U_p(\rho),\qquad
 \delta=\min\left(\frac\beta8,\frac\sigma8,
 \frac\varepsilon{8mU(1+2mM)}\right).
\]
Keep zero exact, approximate other nodes within \(\delta\), and
approximate all coefficients except \(b_{00}\) within
\(\delta/(m+1)\). Set the new anchor coefficient to minus the
sum of the other zeroth coefficients. Every coefficient error is
at most \(\delta\); constants cancel exactly. Nodes remain
distinct inside radius \(\rho\) and
\(|\widetilde b_{jp}|\leq2M\).
Lemma~\ref{lem:perturb} bounds the threshold error by
\(m\delta U+2m^2M\delta U\leq\varepsilon/8\).
Theorem~\ref{thm:ratjets}, with tolerance \(\varepsilon/2\),
approximates the new rational functional. It need not itself arise
from a real circle profile. The total error is less than
\(\varepsilon\), with polynomial encoding throughout.
Approximating and summing the zeroth residues separately gives
the asserted approximation of the real mean \(\mu\).

Finally, a dense rational homogeneous field gives Laurent
\(F,G\) of degree at most \((d+1)/2\), with polynomial bit cost.
Multiplication by that power of \(z\) represents \(F/G\) with
polynomial numerator and denominator degrees at most \(d+1\).
After reflection if necessary, \(G>0\); its denominator has
no unit-circle root. Approximate \(-\mu\) and \(A\) to
\(\gamma/4\) to distinguish the promised cases. For
\(\gamma\geq1\), fixed smaller tolerances suffice.
The strict-certificate criterion also rejects \(a\leq0\);
the equivalence with weak analytic certificates is used only
for \(a>0\), as in Theorem~\ref{thm:decision}.
\end{proof}

Neither theorem makes exact equality polynomial-time, bounds the
output size of a polynomial certificate polynomially, or implements
the classical oracle algorithm. The gain is precision dependence:
increasing the Jackson hierarchy alone gives no corresponding
polynomial bound in \(\log(1/\varepsilon)\).
The primary ingredients remain classical interpolation,
convex-body optimization and root refinement; the reduction
establishes their hypotheses for this certificate threshold.

\section{A semialgebraic region of certificate existence}
The preceding construction is uniform in the coefficients, even when
interior poles collide or cancel. This gives a statement about families
without fixing a certificate degree.

\begin{theorem}[Uniform coefficient region]\label{thm:region}
Fix an odd integer \(d\geq1\). In the coefficient space
\(\R^{2(d+1)}\) of planar homogeneous fields of degree \(d\), let
\(\mathcal R_d\) denote the fields whose angular component has one
strict sign, and let \(\mathcal C_d\subset\mathcal R_d\) denote those
admitting a strict polynomial Lyapunov function of some degree.
Then \(\mathcal C_d\) is open and semialgebraic, with an effectively
constructible first-order real definition over \(\mathbb Q\).
The interpolation matrix in this definition has size at most \(d+2\),
including the normalization node at zero.
\end{theorem}
\begin{proof}
The condition defining \(\mathcal R_d\) is a strict sign condition on
the unit circle, hence semialgebraic and open. Treat \(G>0\); a fixed
reflection treats \(G<0\). Put \(k=(d+1)/2\). The Laurent polynomials
\(F,G\) in \(z=e^{2i\theta}\) have degrees at most \(k\), so
\[
 N(z)=z^k F(z),\qquad D(z)=z^k G(z)
\]
are polynomials of degree at most \(d+1\). Their coefficients depend
\(\mathbb Q[i]\)-linearly on the field coefficients. The polynomial
\(D\) is nonzero and has no root on the unit circle.

There are finitely many degree and root-multiplicity patterns for \(D\).
For each pattern introduce its leading coefficient and distinct complex
roots, impose the coefficient identities for their product, and specify
which roots satisfy \(|\zeta|<1\) and which satisfy \(|\zeta|>1\).
All these conditions are real semialgebraic. Taking the finite union of
patterns handles collisions and degree drops without a continuous choice
of root labels.

The mean \(-a\) is the sum of interior residues of \(N/(zD)\).
On a fixed pattern these residues are rational functions of coefficients
and root variables: repeated poles require only finitely many derivatives.
All denominators are nonzero on the pattern, and auxiliary inverse
variables turn each division into a polynomial equality. This classical
mean computation, also explicit in \cite[Supplement, Theorem~20.3]{korda},
therefore defines \(a\) semialgebraically in the parameters.

Set \(R=N/D+a\). In \eqref{eq:lambda} and \eqref{eq:pick}, use the
interior roots of \(zD\) with their full denominator multiplicities.
This includes zero with multiplicity at least one. Removable poles may
be retained: their unused residue coefficients vanish. The extra jets
do not change the feasible objectives, because every quadrant map
supplies them and Lemma~\ref{lem:pick} still characterizes the full
jet tuples. The total multiplicity is at most
\(\deg D+1\leq d+2\). If \(D(0)=0\), the node zero is counted once
with its multiplicity increased by one.

Both the residue objective and the Pick matrix entries are rational
functions of the root and coefficient variables, polynomial of degree at
most two in the jets. Their denominators, including
\(1-\zeta_j\bar\zeta_\ell\), do not vanish. Thus the anchor condition,
quadrant inequalities and positive semidefiniteness are semialgebraic
jointly in all parameters. Existentially quantify the root and mean
variables, require \(a>0\), and then impose \eqref{eq:qe}, with its
quantifiers over the feasible jets. Any admissible root listing describes
the same profile and threshold. The finite union of these formulas
defines exactly \(\mathcal C_d\); real quantifier elimination gives
the asserted effective semialgebraic description over \(\mathbb Q\).

Finally a strict homogeneous certificate \(P\) for \(X\) has positive
minima of \(P\) and \(-L_XP\) on the circle. The same \(P\) remains
strict under small coefficient changes in \(X\), while the sign of
\(G\) persists. Proposition~\ref{prop:analytic} supplies a homogeneous
certificate at every point of \(\mathcal C_d\), proving openness.
\end{proof}

\begin{corollary}[Finitely many parameter transitions]\label{cor:transitions}
For a polynomial or rational one-real-parameter family of fixed odd
degree, with real algebraic coefficients and a semialgebraic parameter
domain over the real algebraic numbers, assume its rational denominators
do not vanish on that domain. The parameters where the field belongs
to \(\mathcal C_d\) form a finite union of intervals and points with
algebraic finite endpoints. If the domain is open, this set is relatively
open and has no isolated point.
\end{corollary}
\begin{proof}
Take the inverse image of \(\mathcal C_d\) under the semialgebraic
coefficient map. One-dimensional semialgebraic sets have the stated
form; continuity gives the additional openness assertion.
\end{proof}

\begin{corollary}[A computable degree for compact families]\label{cor:compact}
Given a compact semialgebraic coefficient set \(K\), defined over real
algebraic numbers, inclusion \(K\subset\mathcal C_d\) is decidable.
If it holds, an algorithm finds an even degree \(2M\) such that every
\(X\in K\) admits a strict homogeneous certificate of degree \(2M\).
The certificate may depend on \(X\).
\end{corollary}
\begin{proof}
Decide inclusion using Theorem~\ref{thm:region}. A strict homogeneous
certificate at each \(X\in K\) works on a neighborhood. Compactness
gives a finite cover with finitely many certificate degrees. Taking
positive integer powers raises these certificates to their least common
multiple degree and preserves strictness. Hence some even degree works
throughout \(K\).

For \(M=1,2,\ldots\), quantifier elimination decides
\[
 \forall X\in K\ \exists P\in\R[x,y]_{2M}\ \forall(x,y):\quad
 x^2+y^2=1\ \Longrightarrow\ P(x,y)>0\ \text{and}\ -L_XP(x,y)>0,
\]
where \(\R[x,y]_{2M}\) denotes homogeneous forms of degree \(2M\)
and the existential variables are their coefficients. The search
terminates by the finite-cover argument. For empty \(K\) it may return
degree two. This is a common degree, not necessarily a common polynomial,
and no efficiency bound is asserted.
\end{proof}

These are consequences of the finite interpolation representation;
they do not give a new decision procedure for stability itself.

\section{Effective separation and a sufficient degree}
The finite brackets also have a quantitative use independent of
quantifier elimination. If \(a<A(q)\), some \(a_N>a\) excludes
certificates of every degree. If \(a>A(q)\), some \(a_N+E_N<a\)
proves existence. For nonzero separation, it suffices to use
\[
 N=O\!\left(\max\left\{1,
 \sqrt{\frac{\norm{q''}_\infty}{|a-A(q)|}}\right\}\right).
\]
This statement concerns exact bounds on the finite optimum.
A numerical implementation needs certified errors smaller than
the separation; floating solver statuses do not supply them.

\begin{theorem}[Constructive sufficient degree]\label{thm:degree}
Let \(\Delta=a-A(q)>0\) and
\(L=2\sum_{k\geq1}k|q_k|<\infty\).
If \(L=0\), \(P=x^2+y^2\) is a certificate.
Otherwise let
\[
 \tau=\min\{1/2,\Delta/(4L)\},\qquad r=1-\tau,\qquad
 M=\max\left\{1,\left\lceil\frac{8\pi r}
                      {\Delta(1-r^2)}\right\rceil\right\}.
\]
There is a product of \(M\) positive quadratic forms with
normalized slack at least \(\Delta/2\).
When \(\Delta\leq2L\), one may use
\(M\leq\lceil32\pi L/\Delta^2\rceil\).
The Lyapunov polynomial degree is \(2M\).
\end{theorem}
\begin{proof}
Take a diagonal subsequence of finite optimizers. The bounds
\(|m_k|\leq1\), \(a_N\leq\max q\), and closedness of every finite
positive cone give compatible full moments. The classical
positive trigonometric moment representation yields a probability
measure \(\mu\) and a positive measure \(\nu\) of mass \(A\), with
\(\nu_k=i\mu_k-q_k\), \(k\geq1\).
Poisson smoothing at radius \(r\) loses at most
\[
 \norm{q-q_r}_\infty\leq L(1-r)\leq\Delta/4.
\]
The smoothed Herglotz function thus has slack at least \(3\Delta/4\).

Choose a cut of the circle with no atom of \(\mu\), and use its
\(M\) midpoint quantiles with equal weights. Their cumulative
discrepancy is at most \(1/(2M)\), also in the presence of atoms.
Integration by parts and the total variation of the Herglotz kernel
give, uniformly on the closed disk,
\[
 \norm{h_M-h_r}_\infty
 \leq\frac1{2M}\int_0^{2\pi}\frac{2r}{|1-re^{it}|^2}\,dt
 =\frac{2\pi r}{M(1-r^2)}\leq\Delta/4.
\]
Use \eqref{eq:herglotz}. The last bound on \(M\) follows from
\(\tau=\Delta/(4L)\) and
\(r/(1-r^2)\leq1/\tau\).
\end{proof}

No optimality of the exponent or constants is claimed.
Strict binary forms are sums of two squares of real forms, so both
the constructed \(P\) and \(-L_XP\) admit such decompositions.
This is the classical binary positivity fact, not a new general
SOS theorem. A small rational perturbation of the coefficients of
\(P\) preserves both strict inequalities on the unit circle.
Thus for a rational input field a rational-coefficient certificate
exists whenever the strict criterion holds. This does not assert
a two-square decomposition over the rationals.

\section{Two explicit profiles}
\subsection{One harmonic}
If \(q\) is a trigonometric polynomial of degree \(s\), the
functional only uses Taylor coefficients at zero:
\[
 \Lambda_R(f)=\sum_{k=1}^s\overline{q_k}f_k,\qquad f_0=c+i.
\]
Put \(g_k=\sum_{\ell=0}^k f_\ell f_{k-\ell}\).
The finite interpolation condition reduces to
\[
 T_s(\mu)\succeq0,\qquad
 \mu_0=2c,\quad \mu_k=g_k/(2i)\quad(k\geq1).
\]
This is the classical Carath\'eodory--Fej\'er condition for
\(-ig\) to have nonnegative real part; its imaginary constant is
unrestricted.

For \(q=b\cos\phi\), \(b\geq0\), the \(2\times2\) condition is
\[
 |f_1|\leq \frac{2c}{\sqrt{1+c^2}}.
\]
The objective is \((b/2)\Ima f_1-c\), so
\begin{equation}\label{eq:single}
 A(b\cos\phi)
 =\sup_{c>0}\left(\frac{bc}{\sqrt{1+c^2}}-c\right)
 =(b^{2/3}-1)_+^{3/2}.
\end{equation}
The maximum for \(b>1\) occurs at \(c=\sqrt{b^{2/3}-1}\);
for \(b\leq1\) the supremum is zero.
Changing the phase replaces \(b\) by \(|b|\).
Consequently the stable cubic family with \(G=1\) and
\(F=-a+b\cos2\theta\) has a polynomial or analytic weak certificate
exactly when \(|b|<(1+a^{2/3})^{3/2}\).
The family itself is from \cite{liu}.

\subsection{A rational profile with nonconstant angular speed}
Consider
\begin{equation}\label{eq:qexample}
 q(\phi)=\frac{3\cos\phi}{1+(3/5)\cos\phi}+\frac54.
\end{equation}
Writing \(C=25/4\), \(r=1/3\), one has
\(q=C(1-P_{-r})\), where
\(P_{-r}(\phi)=(1-r^2)/|1+r e^{i\phi}|^2\).
Thus
\[
 \Lambda_R(f)=-C\,[f(-r)-f(0)].
\]
Set \(f(0)=c+i\), \(f(-r)=x+iy\), with \(c,x,y>0\).
The two-node Pick condition reduces exactly to
\begin{equation}\label{eq:D}
 D(c,x,y):=18cxy-(c^2-1-x^2+y^2)^2-4(c+xy)^2\geq0.
\end{equation}
Therefore
\begin{equation}\label{eq:Aexample}
 A(q)=\frac{25}{4}-
      \inf_{\substack{c,x,y>0\\D(c,x,y)\geq0}}
      \left(c+\frac{25}{4}y\right).
\end{equation}

\begin{proposition}[An exact algebraic instance]\label{prop:example}
Let
\begin{align*}
 \mathcal R(t)={}&16777216t^{12}-8141930496t^{10}
       +1537904087040t^8\\
 &-104707265987520t^6+2126206939995165t^4\\
 &-15011745808074246t^2+81397225546445341.
\end{align*}
It has a unique root \(t_*\) in \((0,19/4)\), and the profile
\eqref{eq:qexample} satisfies \(A(q)=25/4-t_*\). In particular
\begin{equation}\label{eq:exactinterval}
 1.657217310707<A(q)<1.657217310708.
\end{equation}
The endpoints in \eqref{eq:exactinterval} are exact decimal rationals.
\end{proposition}
\begin{proof}
The feasible rational triple \((c,x,y)=(1,1,3/5)\) has
\(D=94/625>0\) and objective \(19/4<25/4\).
Restrict a minimizing sequence to objective at most \(19/4\).
The equivalent disk condition
\begin{equation}\label{eq:disk}
 \left|(x+iy)^2-\left(c^2-1+\frac52ic\right)\right|
 \leq\frac32c
\end{equation}
bounds \(x\), as \(c,y\) are already bounded.
If \(c\to0\), the squared value tends to \(-1\), forcing
\(x\to0\), \(y\to1\); the objective would tend to \(25/4\),
a contradiction. Hence \(c\) is bounded below.
The imaginary part of \eqref{eq:disk} gives \(2xy\geq c\),
so \(x,y\) also stay bounded away from zero.
A minimum is attained in the positive octant and lies on \(D=0\).

At a regular minimizing boundary point, the multiplier equations,
and at a singular point the vanishing gradient, imply
\begin{equation}\label{eq:critical}
 D=0,\quad D_x=0,\quad 4D_y-25D_c=0,\quad 4t-4c-25y=0.
\end{equation}
Exact rational elimination in the variable order \(x,c,y,t\)
gives the univariate eliminant
\begin{equation}\label{eq:eliminant}
 t^2(4t-25)^2(4t+25)^2(t^2+1)^2\,\mathcal R(t).
\end{equation}
The supplementary program recomputes a graded reverse lexicographic
Gr\"obner basis over \(\mathbb Q\), converts it to lexicographic order,
verifies \eqref{eq:eliminant} and its zero remainder, and performs
exact Sturm root counts. These give one root of \(\mathcal R\)
in \((0,19/4)\), no endpoint root, and one root in
\[
 \left(\frac{4592782689292}{10^{12}},
             \frac{4592782689293}{10^{12}}\right).
\]
Since every minimizer satisfies \eqref{eq:critical} and
\(0<t\leq19/4\), none of the other factors in
\eqref{eq:eliminant} vanishes. Existence of the minimum identifies
it with the unique root \(t_*\), giving \eqref{eq:exactinterval}.
No numerical optimizer is used for global optimality.
\end{proof}

This is realized by a concrete cubic vector field. With
\[
 \rho=x^2+y^2,\quad
 H=\rho+\tfrac35(x^2-y^2),\quad
 K=(-a+\tfrac54)H+3(x^2-y^2),
\]
take \(\dot x=Kx-Hy\), \(\dot y=Ky+Hx\).
Then \(G=1+(3/5)\cos2\theta>0\) and
\(F=(-a+5/4)G+3\cos2\theta\).
It is globally asymptotically stable for every \(a>0\);
polynomial and local analytic weak certificates exist precisely
for \(a>A(q)\). The algebraic boundary value is therefore an
explicit equality input to the finite decision theorem.

\subsection{An exact minimum degree}
At \(a=2\), the preceding cubic has a strict integer-coefficient
certificate, while finite rational identities exclude every smaller degree.

\begin{proposition}\label{prop:min8}
For the cubic above at \(a=2\), the minimum degree of a polynomial
weak Lyapunov function is eight, even allowing nonhomogeneous polynomials
and inequalities only near the origin, with \(V(0)=0\).
A homogeneous degree-eight strict certificate exists globally.
\end{proposition}
\begin{proof}
Define
\begin{align*}
 P(x,y)={}&216x^8-2260x^7y+12820x^6y^2-45066x^5y^3
              +99726x^4y^4\\
          &-137242x^3y^5+113341x^2y^6-51276xy^7+9784y^8.
\end{align*}
The descending \(x\)-coefficients of \(5(-L_XP)\), of degree ten, are
\[
\begin{gathered}
2528,\ -28576,\ 65016,\ 322736,\ -1562236,\ 1039608,\\
5502688,\ -13799076,\ 13740768,\ -6471612,\ 1188936.
\end{gathered}
\]
Exact signed Euclidean Sturm chains for \(P(t,1)\) and
\(-L_XP(t,1)\) have variations \((4,4)\) and \((5,5)\) at
minus and plus infinity, respectively. Neither has a real root.
Their values at \((0,1)\) are \(9784\) and \(1188936/5\),
and at the omitted projective point \((1,0)\) they are
\(216\) and \(2528/5\). Homogeneity proves both strict positivities.

For the lower bound write, at degree \(2M\),
\[
 p(\phi)=p_0+2\Rea\sum_{k=1}^M p_k e^{ik\phi},\qquad
 d(\phi)=-MF(\phi)p(\phi)-G(\phi)p'(\phi),
\]
where \(\phi=2\theta\), \(F=-3/4+(51/20)\cos\phi\), and
\(G=1+(3/5)\cos\phi\). On the circle \(-L_XP=2d\).
Let \(v_M=(1,2\cos\phi,-2\sin\phi,\ldots,
2\cos M\phi,-2\sin M\phi)^T\).
For a real rational \(t\), set
\(\cos\phi=(1-t^2)/(1+t^2)\), \(\sin\phi=2t/(1+t^2)\), and
\[
 b_M(t)=-MF(\phi)v_M(\phi)-G(\phi)v_M'(\phi).
\]
Thus \(b_M(t)\) is a rational column, and its pairing with the
real Fourier coefficients gives \(d(\phi)\).
For the nodes in the following table form
\(B_M=[b_M(t_1)\ \cdots\ b_M(t_{2M+1})]\), and define
\(w=-B_M^{-1}(1,0,\ldots,0)^T\).
These formulas specify the weights exactly over \(\mathbb Q\).
\begin{center}
\begin{tabular}{clcc}
\toprule
\(M\)& Nodes \(t_j\)&\(\operatorname{sign}\det B_M\)&\(\min_j w_j>\)\\
\midrule
1&\(-3/10,\ 17/20,\ 9/10\)&\(1\)&1\\
2&\(-7/20,\ 9/20,\ 1/2,\ 23/20,\ 6/5\)&\(-1\)&2\\
3&\(-7/20,\ 1/4,\ 3/10,\ 7/10,\ 3/4,\ 13/10,\ 27/20\)&\(1\)&6\\
\bottomrule
\end{tabular}
\end{center}
Direct rational arithmetic gives the nonzero determinants and positive
weight bounds displayed; the supplementary verifier supplies the full
rational weights and checks every coefficient of
\[
 \sum_j w_j d(\phi_j)=-p_0.
\]
For a nonzero nonnegative \(p\), its mean \(p_0\) is positive.
The identity is therefore incompatible with \(d\geq0\), excluding
degrees two, four and six over arbitrary real coefficients.

Finally, the first nonzero homogeneous term of any local polynomial
weak certificate with \(V(0)=0\) is nonnegative, nonzero, of positive
even degree, and has nonpositive Lie derivative. If its total degree
were below eight, one of the three excluded cases would occur.
The exact verifier uses no optimizer; numerical searches only discovered
the fixed coefficients and support nodes.
\end{proof}

This is a minimum-degree result for this instance. It does not give
a general minimum-degree law or a new positivity technique.

\section{A robust obstruction when angular velocity changes sign}
The preceding decision and complexity results require nonvanishing
angular velocity. A sufficient obstruction extends to angular zeros
and yields an explicit open set of stable systems outside that class.
This is not a complete characterization with angular zeros.

\begin{lemma}[A mean obstruction with boundary factors]\label{lem:zeros}
For a homogeneous field of odd degree, let \(F,G\) be its real
trigonometric polar profiles in \(\phi=2\theta\),
with \(G\not\equiv0\), and write
\[
 G(\phi)=\sum_{k=-s}^s\gamma_k e^{ik\phi},\qquad
 T_G(\alpha)=2\Ima\sum_{k=1}^s\overline{\gamma_k}\alpha^k.
\]
If
\begin{equation}\label{eq:zeroobstruction}
 \av{F}+\min_{|\alpha|\leq1}T_G(\alpha)>0,
\end{equation}
then the homogeneous field has no nonzero nonnegative homogeneous
polynomial of positive degree with nonpositive Lie derivative.
It consequently has no positive local real-analytic weak Lyapunov function.
\end{lemma}
\begin{proof}
A nonzero nonnegative homogeneous polynomial has degree \(2M>0\)
and circle restriction
\[
 p(\phi)=C\prod_{j=1}^M|1-\alpha_j e^{i\phi}|^2,
 \qquad C>0,\quad |\alpha_j|\leq1.
\]
Zero parameters may complete the factor list. Where \(p>0\), the
Lie inequality is
\[
 F+G\Ima h\leq0,\qquad
 h(z)=\frac1M\sum_{j=1}^M\frac{1+\alpha_jz}{1-\alpha_jz}.
\]
Indeed \(\dot\phi=2r^{d-1}G\) and
\(p'/(Mp)=\Ima h\). A zero of \(p\) at which \(G\ne0\)
is impossible: on a short interval with constant sign of \(G\),
the inequality \(Gp'+MFp\leq0\), an integrating factor, and
\(p\geq0\) propagate that zero to a one-sided interval. The
analytic trigonometric polynomial would vanish identically.

For \(|\alpha|<1\), Fourier expansion gives
\[
 \av{G\Ima\frac{1+\alpha e^{i\phi}}{1-\alpha e^{i\phi}}}
 =T_G(\alpha).
\]
For a boundary factor \(|\alpha|=1\), its pole occurs at a zero
of \(G\). In a local coordinate \(t\) about this point,
\(G(t)=O(t)\), so its product with \(\cot(t/2)\) extends
continuously. The radial regularization has imaginary part
\[
 \frac{2r\sin t}{(1-r)^2+2r(1-\cos t)}.
\]
For \(|t|\leq\pi\), the inequalities
\(|G(t)|\leq L|t|\) and
\(1-\cos t\geq2t^2/\pi^2\) bound the product uniformly by
\(L\pi^2/2\). Away from its pole it converges uniformly.
Dominated convergence proves the same mean identity for boundary
factors. Averaging the Lie inequality therefore gives
\[
 0\geq\av{F}+\frac1M\sum_jT_G(\alpha_j)
 \geq\av{F}+\min_{|\alpha|\leq1}T_G(\alpha)>0,
\]
a contradiction. The first nonzero homogeneous Taylor term of a
positive analytic weak certificate is nonnegative, of positive even
degree, and has nonpositive Lie derivative. The same contradiction
excludes it; strict positivity of that leading term is not assumed.
\end{proof}

\begin{theorem}[A full coefficient box with a common rational certificate]
\label{thm:box}
Consider the quintic field
\[
 f_0=-x^5+16x^3y^2-3xy^4,\qquad
 g_0=x^4y+20x^2y^3-y^5.
\]
Independently perturb each of the twelve real coefficients of these
two homogeneous quintics by at most \(1/240\). Every resulting
field is globally asymptotically stable, its angular velocity changes
sign, and it admits no positive local analytic weak Lyapunov function.
The same function
\begin{equation}\label{eq:commonrational}
 V(x,y)=\frac{(11x^2+y^2)^{11}}{(x^2+y^2)^{10}}
 \quad((x,y)\ne0),\qquad V(0)=0,
\end{equation}
is a proper positive \(C^1\) certificate for the entire box, satisfying
\(L_XV\leq-r^4V/5\).
\end{theorem}
\begin{proof}
The central polar profiles are
\[
 F_0=-1+5\sin^2\phi,\qquad G_0=\sin\phi.
\]
Each perturbing Cartesian component has at most six terms, each
bounded by \(1/240\) on the unit circle. Hence the changes in both
\(F\) and \(G\) have uniform norm at most
\(\epsilon=1/20\). The profiles have trigonometric degree at most
three in \(\phi\). Consequently
\[
 |\av{F}-3/2|\leq\epsilon,\qquad
 \norm{T_G-T_{G_0}}_{\infty,\overline\D}\leq6\epsilon.
\]
Since \(T_{G_0}(\alpha)=\Rea\alpha\), the obstruction margin in
\eqref{eq:zeroobstruction} is at least
\(1/2-7\epsilon=3/20>0\). Also
\(G(\pi/2)\geq19/20\) and
\(G(3\pi/2)\leq-19/20\), proving sign change.

For stability set \(v=(\partial_\theta\log V)/2\). Directly,
\[
 v=-\frac{55\sin\phi}{6+5\cos\phi},\qquad
 F_0+G_0v=-1-\frac{25\sin^2\phi(1-\cos\phi)}{6+5\cos\phi}
 \leq-1.
\]
The identity
\((6+5t)^2-11(1-t^2)=(5+6t)^2\) implies
\(|v|\leq5\sqrt{11}<17\). Thus, for every field in the box,
\[
 F+Gv\leq-1+18\epsilon=-\frac1{10},\qquad
 L_XV=2r^4V(F+Gv)\leq-\frac15r^4V.
\]
The function satisfies \(r^2\leq V\leq11^{11}r^2\) and
\(\nabla V=O(r)\), so its extension is \(C^1\) and proper.
Its compact sublevel sets ensure forward completeness; the strict
Lie inequality gives global asymptotic stability.
\end{proof}

\begin{corollary}[Higher-order local perturbations]\label{cor:higher}
If a locally \(C^1\) vector field satisfies \(X=X_5+O(r^6)\),
with no terms of degree below five and
\(X_5\) in the box of Theorem~\ref{thm:box}, the origin is locally
asymptotically stable and has no positive local analytic weak
Lyapunov function.
\end{corollary}
\begin{proof}
The extra contribution to \(L_XV\) is \(O(r^7)\), whereas the
negative term is at most \(-r^6/5\). The first nonzero Taylor
term of a putative analytic certificate is governed by \(X_5\)
and is excluded by Lemma~\ref{lem:zeros}. This argument is local;
arbitrary higher-order perturbations need not preserve global stability.
\end{proof}

Liu--Fitzsimmons already establish robust counterexamples within
their two-parameter rotating family \cite[Section~IV.D]{liu},
using the factorization and Taylor methods that underlie our lemma.
Ahmadi--El Khadir's rational converse theorem
\cite[Theorem~3.1]{ahmadi-khadir} already guarantees rational
certificates for asymptotically stable homogeneous fields.
The specific additional conclusions here are the boundary-factor
obstruction and a full twelve-dimensional rational coefficient box
with angular sign changes and an explicit \emph{common} certificate.
The rational expression is singular at the origin before extension;
it is not an analytic certificate there. Smooth nonanalytic
Lyapunov functions are not excluded.

\section{Verification, scope and remaining questions}
The principal equivalences are analytic proofs.
The rational threshold additionally uses exact symbolic elimination
and Sturm counting; its program, equations, output and rational
interval are supplied for reproduction. Proposition~\ref{prop:min8}
uses separate exact positivity checks and rational evaluation identities,
including a verification that does not call an optimizer.
Separate exact checks verify Fourier signs, Toeplitz quadratic-form
identities, atomic quantile discrepancy, the confluent Pick kernel
including a singular Gram example, determinant identities and residues.
The coefficient-region proof additionally accounts for pole collisions,
removable poles and degree drops. Its quantifier elimination is theoretical.
The complexity results explicitly invoke classical weak convex optimization
and certified root refinement; neither general algorithm is implemented here.
Additional exact checks cover the rational PSD oracle, confluent Hermite
perturbations, irrational simple and repeated residue poles, functional
perturbation bounds on Cayley kernels, and the common rational quintic
certificate. These finite checks verify formulas and signs, not the full
analytic or bit-complexity proofs. Reviewer checks are recorded separately.

Floating computations are retained as diagnostics. The hierarchy
experiment has fifteen solves, nine labelled \texttt{optimal} and
six \texttt{optimal\_inaccurate}. Eight further primal/dual comparisons
include three inaccurate primal and two inaccurate dual solves.
Small negative numerical duality gaps are not certified bounds.
A six-start numerical search with nonlinear constraints for the rational
example produced five
infeasible failures and one apparent local success; all are retained.
Proposition~\ref{prop:example} uses the separate exact elimination
argument, not that local success.

AI tools assisted the derivations, computation and writing.
Proof dossiers were examined by the same involved model reviewer
requested as GPT-6 Astra with extra-high effort. This is not blind
review, an additional independent vote, human peer review or formal
verification; runtime identity is not independently attested.
Any exact-manuscript assessment is a separate record.
No Lean proof is supplied.

The literature comparison is targeted, not exhaustive.
The original stability counterexample is credited to \cite{liu};
the kernels, interpolation and quantifier elimination are classical.
Worldwide priority of this connection is unconfirmed.
The all-degree characterization remains restricted to nonvanishing
angular velocity. Theorem~\ref{thm:box} supplies a sufficient obstruction
with angular zeros, and Corollary~\ref{cor:higher} treats only the stated
local higher-order perturbations. No general nonhomogeneous or
higher-dimensional classification follows. The polynomial-time claims
concern additive values and supplied gap promises in dense encoding;
they do not decide equality efficiently or produce a polynomial-sized
certificate. Practical certified solvers, optimal minimum-degree laws,
and a complete certificate criterion with angular zeros remain open.

\small
\begin{thebibliography}{99}
\bibitem{liu}
J.~Liu and M.~Fitzsimmons,
\emph{A globally asymptotically stable planar homogeneous polynomial
vector field with no polynomial Lyapunov function},
arXiv:2607.16171v1, 17 July 2026.
\url{https://arxiv.org/abs/2607.16171v1}.
\bibitem{ahmadi-parrilo}
A.~A. Ahmadi and P.~A. Parrilo,
\emph{Stability of polynomial differential equations:
complexity and converse Lyapunov questions}, arXiv:1308.6833.
\url{https://web.mit.edu/~a_a_a/Public/Publications/Selected_Papers/dTAC.pdf}.
\bibitem{ahmadi-khadir}
A.~A. Ahmadi and B.~El Khadir,
\emph{On algebraic proofs of stability for homogeneous vector fields},
IEEE Trans. Automat. Control \textbf{65} (2020), 325--332.
\url{https://optimization-online.org/wp-content/uploads/2018/03/6502.pdf}.
\bibitem{jackson}
E.~de Klerk, R.~Hess and M.~Laurent,
\emph{Improved convergence rates for Lasserre-type hierarchies of
upper bounds for box-constrained polynomial optimization},
SIAM J. Optim. \textbf{27}(1) (2017), 347--367.
\url{https://doi.org/10.1137/16M1065264}.
\bibitem{clark}
D.~N. Clark, \emph{On matrices associated with generalized
interpolation problems}, Pacific J. Math. \textbf{27}(2)
(1968), 241--253.
\url{https://msp.org/pjm/1968/27-2/pjm-v27-n2-p04-p.pdf}.
\bibitem{tarski}
A.~Tarski, \emph{A decision method for elementary algebra and geometry},
second edition, University of California Press, 1951.
\url{https://doi.org/10.1525/9780520348097}.
\bibitem{fejer}
M.~A. Dritschel and J.~Rovnyak,
\emph{The operator Fej\'er--Riesz theorem}, arXiv:0903.3639.
\url{https://arxiv.org/abs/0903.3639}.
\bibitem{tripathi}
S.~Tripathi, F.~Ilhan, A.~Gokcen, M.~Kudeyt, S.~Sahin and O.~Karabacak,
\emph{Construction of Lyapunov certificates for oscillatory systems},
preprint, 26 February 2026.
\url{https://doi.org/10.20944/preprints202602.1723.v1}.
\bibitem{korda}
M.~Korda, \emph{Lyapunov stability of polynomial vector fields is
undecidable}, arXiv:2609.22058v2, 25 September 2026, and
\emph{Supplementary material}, Part~III, pp.~51--59.
\url{https://arxiv.org/abs/2609.22058v2};
\url{https://homepages.laas.fr/mkorda/arnold_supplement.pdf}.
\bibitem{gls}
M.~Gr\"otschel, L.~Lov\'asz and A.~Schrijver,
\emph{Geometric Algorithms and Combinatorial Optimization},
Springer, 1988, Corollary~4.3.12 and definitions in Section~2.1.
\url{https://www.zib.de/userpage/groetschel/pubnew/paper/groetschellovaszschrijver1988.pdf}.
\bibitem{msw}
K.~Mehlhorn, M.~Sagraloff and P.~Wang,
\emph{From approximate factorization to root isolation with application
to cylindrical algebraic decomposition}, arXiv:1301.4870v2,
23 January 2014, Theorems~3--4.
\url{https://arxiv.org/abs/1301.4870v2}.
\end{thebibliography}
\end{document}
